% META: {"id": "1NSI-WEB-EX-001", "chapitre": "1NSI-WEB-IHM", "type_objet": "exercice", "capacites": ["P-IHM-01A", "P-IHM-01B"], "parcours": 1, "competences": ["raisonner"], "duree_min": 10, "mode_creation": "ex_nihilo", "sources_inspiration": [], "fichier_tex": "chapitres/1NSI-WEB-IHM/exercices/1NSI-WEB-EX-001.tex", "corrige_tex": "chapitres/1NSI-WEB-IHM/corriges/1NSI-WEB-CO-001.tex", "status": "generated"}
\begin{exercice}{1NSI-WEB-EX-001}{1}{10}
On donne l'extrait de code HTML suivant :
\begin{console}
<form id="formulaire_avis">
    <input type="text" id="champ_pseudo" name="pseudo">
    <select id="liste_note" name="note">
        <option value="5">5 etoiles</option>
        <option value="1">1 etoile</option>
    </select>
    <button id="bouton_envoyer" type="submit">Envoyer</button>
</form>
\end{console}
\begin{enumerate}
  \item Identifier les \textbf{trois} composants graphiques d'interaction présents dans ce formulaire (hors le formulaire lui-même).
  \item Pour chacun, indiquer un événement plausible qu'une fonction JavaScript pourrait traiter (par exemple : saisie de texte, changement de sélection, clic).
  \item Le code HTML seul suffit-il à faire réagir la page lorsqu'on clique sur le bouton ? Justifier.
\end{enumerate}
\end{exercice}
